\scanpage{108}
% Archive pencil foliation: 107.
motifs sur $\underline{C}$ prend ses valeurs en fait dans la catégorie des structures de Hodge (une polarisation d'une varéiét projective lisse $V$ définit canoniquement une polarisation de la structure de Hodge associée sur la cohomologie de Betti-Hodge). NB On n'a aucune idée sur ce que pourrait être l'image essentielle du foncteur précédent, par exemple s'il y a lieu d'esp\unsure{è|\~{e}}rer qu'on trouve toutes les structures de Hodge (doncune équivalence de catégories : \begin{tikzcd}[cramped]\text{motifs sur }\underline{C}\arrow[r]&\text{structures de Hodge}\end{tikzcd}); cela semble peu probable, mais on n'a aucune indicationsérieuse dans un sens ou l'autre.

La catégorie des \deleted{pré} structures de Hodge est semi-simple. \deleted{En fait,} Si $G$ est le groupe de Galois d'une \inserted{\rlap{\handwritten{\begin{tikzpicture}[baseline=-.2ex,x=1em,y=1em]
\draw (-.35,.2) .. controls (.1,.15) and (.8,-.05) .. (1.25,-.5) .. controls (1.8,.2) and (2.2,.65) .. (2.5,1.15);
\end{tikzpicture}}}$\otimes$-}catégorie \inserted{\rlap{\handwritten{\begin{tikzpicture}[baseline=-.2ex,x=1em,y=1em]
\draw (-.35,.2) .. controls (0,.35) and (.25,-.2) .. (.45,-.8) .. controls (.75,-.2) and (1.1,.4) .. (1.7,1.2);
\end{tikzpicture}}}$\underline{C}$} de pré-structures de Hodge, à engendrement fini si on veut (pour simplifier), alors on peut expliciter en termes du groupe de Galois associé et de s\unsure{\deleted{e}a|\deleted{o}a|\deleted{a}e|\deleted{a}o}\deleted{\unsure{s|n}} structure\deleted{s} $i_1\deleted{,i_2}$ ci-dessus la condition \deleted{que} pour que les objets de $\underline{C}$ soient en fait des structures de Hodge. Ceci est un exercice plaisant et délectable, qui devrait figurer dans un travail systématique sur les $\otimes$-catégories, dans le chapitre des exemples. On trouve des restrictions très sérieuses sur le groupe $G$ \deleted{(plus} muni de $i_1$ (en plus du fait que $G$ soit réductif).

4) Je laisse le soin à Saavedra de déterminer quelle structure supplémentaire on obtient sur la structure de Hodge ``complexe'' associée à une variété projective \inserted{\rlap{\handwritten{\begin{tikzpicture}[baseline=-.2ex,x=1em,y=1em]
\draw (.1,-.7) .. controls (.2,-.25) and (.5,.1) .. (1,.15) .. controls (2.3,.25) and (3.3,.25) .. (3.5,.5) -- (3.55,.9);
\end{tikzpicture}}}lisse} complexe $X$, lorsqu'on se donne cette dernière comme déduite d'une variété project\rlap{\raisebox{1.1ex}{\tikz[x=1em,y=1em]{\draw (0,0)--(0,.12) (0,.2)--(.01,.33) (.01,.43)--(.02,.49);}}}ive réelle $X_{\underline{R}}=X_o$.On trouve une notion de ``structure de Hodge\ \ réelle'', donnant naissance à une $\otimes$-catégorie correspondante. Dans le groupe de Galois motivique de celui-ci, en plus de la structure $i_1$, on trouve un élément $f_{\handwritten{\infty}}$ de $G(\underline{Q})$, \deleted{de degré} d'ordre 2 (jouant le rôle d'un ``élément de \unsure{\deleted{f}F|\deleted{F}f}robenius à l'infini''), qui correspond \margin{\tikz[x=1em,y=1em]{\draw (0,0) .. controls (.08,.02) and (.15,.13) .. (.2,.25);}}à l'automorphisme du foncteur de Betti \begin{tikzcd}[cramped]X_o\arrow[r,mapsto]&H^{\handwritten{*}}(X_o(\underline{C}),\underline{Q})\end{tikzcd} déduit de l'homéomorphisme \begin{tikzcd}[cramped]x\arrow[r,mapsto]&\overset{\handwritten{\rule{.55em}{.35pt}}}{x}\end{tikzcd} de $X_o(\underline{C})$. Il faut expliciter les relations entre cet élément et $i_1$, $i_2$ !

